\section{Methodology: GraviMerge}
\label{sec:method}

We present \textbf{GraviMerge}, a physics-informed test-time model ensembling framework based on second-order Newtonian gravitational dynamics on spherical manifolds. Instead of modeling the activation flow through deep layers using static distances or linear chemical decay, GraviMerge establishes an auxiliary physical cosmology inside the high-dimensional latent space. Specifically, we decouple the actual neural representation propagation from the physics-informed routing dynamics. We map the routing controller to an auxiliary virtual spacecraft coordinate probe traveling through the latent space, while mapping task-specific expert adapters to celestial bodies of varying masses. By introducing second-order momentum, geodesic manifold integration, and parallel transport of the spacecraft probe, GraviMerge generates highly stable, continuous, and robust ensembling weights that govern the model merging process at each layer.

---

\vspace{-5pt}
\subsection{Physical Analogy and Coordinate Space}
Let the backbone network consist of $L$ sequential layers with an intermediate activation dimension of $D$. We assume there are $K$ pre-trained task-specific expert adapters (such as LoRA adapters \cite{hu21lora}) that can be dynamically merged at test-time. Let $\boldsymbol{\mu}_k^{(3)} \in \mathbb{R}^D$ represent the pre-computed centroid for task $k \in \{1, \dots, K\}$, extracted at Layer 3 (the boundary of our early shared extractor block).

While we present static centroids here for our sandbox, GraviMerge naturally scales to layer-specific centroids $\boldsymbol{\mu}_k^{(l)}$. Under this formulation, the celestial attractors vary dynamically across depth, enabling GraviMerge to track representational drift across layer manifolds of deep transformers, substituting $\boldsymbol{\mu}_k^{(3)}$ with $\boldsymbol{\mu}_k^{(l)}$ seamlessly.

We establish the following mappings between deep learning representations and orbital mechanics:
\begin{itemize}
    \item \textbf{Coordinate Space: } The representation space is bounded to the unit hypersphere $\mathbb{S}^{D-1}$. This aligns with modern cosine similarity routing paradigms and prevents unbounded feature explosion.
    \item \textbf{Spacecraft Coordinate Probe: } An auxiliary state vector $\mathbf{h}_{\text{sc}}^{(l-1)} \in \mathbb{S}^{D-1}$ tracks the coordinates of a virtual spacecraft traveling through the latent gravity field at layer step $l$.
    \item \textbf{Task Centroids (Celestial Stars): } The task centroids $\boldsymbol{\mu}_k^{(3)}$ act as stationary celestial bodies (stars or massive planets) located at fixed coordinates on $\mathbb{S}^{D-1}$, exerting gravitational forces on the spacecraft probe.
    \item \textbf{Layer Depth as Virtual Time: } The sequential flow of representations through layer index $l \in [4, L]$ is formulated as a discrete-time orbital integration, where each layer step represents a virtual time increment $\Delta t$.
    \item \textbf{Probe Momentum: } The spacecraft maintains a virtual velocity vector $\mathbf{v}^{(l)} \in \mathbb{R}^D$, providing the physical inertia required to smooth out sudden ensembling weight changes.
\end{itemize}

---

\vspace{-5pt}
\subsection{Arrhenius Mass Activation (AMA)}
When an input query sample is processed by the model, it first flows through the early shared layers (Layers 1--3). The output activation $\mathbf{h}^{(3)}$ is used to initialize the spacecraft coordinates $\mathbf{h}_{\text{sc}}^{(3)} = \mathbf{h}^{(3)} / \|\mathbf{h}^{(3)}\|_2$ and probe the input's task affinity. Rather than using static gravitational masses for the stars, we compute dynamic gravitational masses $M_k \ge 0$ at test-time. This dynamic mass allocation ensures that stars representing highly relevant tasks are assigned massive gravitational charges, while unrelated tasks behave as light-mass dwarf stars.

We define the dynamic mass $M_k$ of task attractor $k$ using an Arrhenius/Gibbs factor scaled by the angular similarity:
\vspace{-3pt}
\begin{equation}
    M_k = \exp\left( \frac{\cos(\mathbf{h}_{\text{sc}}^{(3)}, \boldsymbol{\mu}_k^{(3)}) - \max_j \cos(\mathbf{h}_{\text{sc}}^{(3)}, \boldsymbol{\mu}_j^{(3)})}{\tau_{\text{grav}}} \right)
    \label{eq:ama}
\end{equation}
\vspace{-3pt}
where $\tau_{\text{grav}} > 0$ is the routing temperature. In practice, subtracting the maximum similarity inside the exponent is a crucial numerical stabilizer that prevents floating-point overflow and bounds the active masses to $M_k \in (0, 1]$, with the most relevant task star acting as a normalized unit attractor. A small $\tau_{\text{grav}}$ exponentially amplifies the relative gravitational mass of the most aligned task, creating a dominant attractor in the latent space.

---

\vspace{-5pt}
\subsection{Multi-Body Softened Gravitational Forces}
At each adapted layer $l \in [4, L]$, the spacecraft probe is positioned at $\mathbf{h}_{\text{sc}}^{(l-1)} \in \mathbb{S}^{D-1}$. To calculate the gravitational forces, we first compute the distance between the spacecraft and each task star. Because coordinates are normalized, we formulate the distance $r_{k, l-1}$ as the exact Euclidean distance on the unit hypersphere:
\vspace{-3pt}
\begin{equation}
    r_{k, l-1} = \left\| \mathbf{h}_{\text{sc}}^{(l-1)} - \boldsymbol{\mu}_k^{(3)} \right\|_2 = \sqrt{2 - 2\cos(\mathbf{h}_{\text{sc}}^{(l-1)}, \boldsymbol{\mu}_k^{(3)})}
    \label{eq:distance}
\end{equation}
\vspace{-3pt}

Newton's law of gravitation defines force as inversely proportional to $r_k^2$. If the spacecraft passes directly through a task centroid ($r_k \to 0$), the force would experience a numerical division-by-zero singularity, leading to chaotic trajectories. To guarantee numerical stability, we incorporate a softened inverse-square gravitational force vector $\mathbf{F}_k^{(l)} \in \mathbb{R}^D$ defined as:
\vspace{-3pt}
\begin{equation}
    \mathbf{F}_k^{(l)} = G \frac{M_k}{r_{k, l-1}^2 + \epsilon^2} \hat{\mathbf{u}}_k^{(l-1)}
    \label{eq:force}
\end{equation}
\vspace{-3pt}
where $G > 0$ is the gravitational constant, $\epsilon > 0$ is the softening parameter, and $\hat{\mathbf{u}}_k^{(l-1)}$ is the unit attractive direction vector pointing from the spacecraft to the star:
\vspace{-3pt}
\begin{equation}
    \hat{\mathbf{u}}_k^{(l-1)} = \frac{\boldsymbol{\mu}_k^{(3)} - \mathbf{h}_{\text{sc}}^{(l-1)}}{\left\| \boldsymbol{\mu}_k^{(3)} - \mathbf{h}_{\text{sc}}^{(l-1)} \right\|_2}
    \label{eq:unit_vector}
\end{equation}
\vspace{-3pt}

This softened force magnitude decays as $1/r^2$ at large distances ($r \gg \epsilon$), and smoothly approaches its peak value $G M_k / \epsilon^2$ at the center ($r \to 0$). Under standard classical mechanics, physical force is defined as the negative gradient of potential energy, $\mathbf{F}_k = -\nabla_{\mathbf{h}_{\text{sc}}} \Phi(r)$. Thus, our attractive gravitational force matches the negative gradient of a positive \textit{Arctangent Potential}:
\vspace{-3pt}
\begin{equation}
    \Phi(r) = \frac{G M_k}{\epsilon} \arctan\left(\frac{r}{\epsilon}\right)
    \label{eq:arctangent_potential}
\end{equation}
\vspace{-3pt}
Indeed, since the gradient of the distance is $\nabla_{\mathbf{h}_{\text{sc}}} r = -\hat{\mathbf{u}}_k$, we have $-\nabla_{\mathbf{h}_{\text{sc}}} \Phi(r) = - \frac{d\Phi}{dr} (-\hat{\mathbf{u}}_k) = \frac{G M_k}{r^2 + \epsilon^2} \hat{\mathbf{u}}_k$, yielding exactly the attractive force vector in Equation~\ref{eq:force}. Under a traditional Plummer potential, the gradient unphysically drops to zero as $r \to 0$ due to spatial symmetry. In our weight-blending paradigm, closer proximity to an expert star must increase its influence. Our Arctangent potential ensures the force magnitude smoothly approaches its peak value at the centroid, providing a well-behaved maximum routing influence without numerical singularities or unphysical force cancellations.

---

\vspace{-5pt}
\subsection{Geodesic Trajectory Integration (GTI)}
Under Newton's second law, the net force accelerating the spacecraft probe is the vector sum of all gravitational forces: $\mathbf{a}^{(l)} = \sum_{k=1}^K \mathbf{F}_k^{(l)}$. To track the smooth, continuous trajectory, the probe maintains a velocity vector $\mathbf{v}^{(l)} \in \mathbb{R}^D$. At the boundary layer, the spacecraft starts from rest, initializing $\mathbf{v}^{(3)} = \mathbf{0} \in \mathbb{R}^D$.

To enforce mathematically rigorous, manifold-consistent dynamics on the sphere $\mathbb{S}^{D-1}$, we project the flat Euclidean acceleration and velocity updates onto the local tangent space of the sphere. This prevents the velocity state vector from accumulating unphysical radial components, resolving potential state saturation and numerical drift.

At each layer step $l \in [4, L]$, we project the net acceleration and update velocity on the tangent plane at position $\mathbf{h}_{\text{sc}}^{(l-1)}$ to preserve spherical constraints:
\vspace{-3pt}
\begin{align}
    \mathbf{a}_{\text{tangent}}^{(l)} &= \mathbf{a}^{(l)} - \left( \mathbf{a}^{(l)} \cdot \mathbf{h}_{\text{sc}}^{(l-1)} \right) \mathbf{h}_{\text{sc}}^{(l-1)} \label{eq:acc_projection} \\
    \mathbf{v}_{\text{tangent}}^{(l)} &= \gamma_{\text{drag}} \mathbf{v}^{(l-1)} + \mathbf{a}_{\text{tangent}}^{(l)} \Delta t \label{eq:velocity_update} \\
    \mathbf{v}^{(l)} &= \mathbf{v}_{\text{tangent}}^{(l)} - \left( \mathbf{v}_{\text{tangent}}^{(l)} \cdot \mathbf{h}_{\text{sc}}^{(l-1)} \right) \mathbf{h}_{\text{sc}}^{(l-1)} \label{eq:velocity_projection}
\end{align}
\vspace{-3pt}
While in infinite-precision arithmetic the projection in Equation~\ref{eq:velocity_projection} is mathematically redundant (as the sum of tangent vectors $\mathbf{v}^{(l-1)}$ and $\mathbf{a}_{\text{tangent}}^{(l)}$ is strictly tangent), under finite 32-bit (or 16-bit) floating-point precision, cumulative addition can introduce subtle numerical drift. This projection acts as a crucial numerical stabilizer, ensuring strict, high-fidelity tangent-space adherence and preventing coordinate scale accumulation.

To update the coordinate position tracker state on the curved manifold $\mathbb{S}^{D-1}$, we utilize the exact spherical \textit{Exponential Map} (geodesic step):
\vspace{-3pt}
\begin{equation}
    \mathbf{h}_{\text{sc}}^{(l)} = \cos(\|\mathbf{v}^{(l)}\|_2 \Delta t) \mathbf{h}_{\text{sc}}^{(l-1)} + \sin(\|\mathbf{v}^{(l)}\|_2 \Delta t) \frac{\mathbf{v}^{(l)}}{\|\mathbf{v}^{(l)}\|_2}
    \label{eq:exponential_map}
\end{equation}
\vspace{-3pt}
Because tangent spaces at different points on curved manifolds are not parallel, we must also parallel transport the velocity vector from the tangent space at $\mathbf{h}_{\text{sc}}^{(l-1)}$ to the tangent space at $\mathbf{h}_{\text{sc}}^{(l)}$. We implement the exact closed-form spherical \textit{Parallel Transport} of $\mathbf{v}^{(l)}$ along the geodesic connecting the two coordinates:
\vspace{-3pt}
\begin{equation}
    \mathbf{v}^{(l)}_{\text{transported}} = \mathbf{v}^{(l)} - \frac{\mathbf{h}_{\text{sc}}^{(l-1)} + \mathbf{h}_{\text{sc}}^{(l)}}{1 + \mathbf{h}_{\text{sc}}^{(l-1)} \cdot \mathbf{h}_{\text{sc}}^{(l)}} \left( \mathbf{v}^{(l)} \cdot \mathbf{h}_{\text{sc}}^{(l)} \right)
    \label{eq:parallel_transport}
\end{equation}
\vspace{-3pt}
This transported velocity vector serves as the velocity state input $\mathbf{v}^{(l)}$ for the subsequent layer step. This rigorous formulation preserves the geometric invariants of second-order spherical motion, preventing numerical drift and coordinate saturation, and ensuring absolute numerical stability.

\textbf{Closed-Loop Feedback Coupling (Coupled GraviMerge): } To allow the spacecraft probe to dynamically respond to unexpected representational shifts inside deep layers (resolving the decoupling critique), we introduce a closed-loop alternative called \textit{Coupled GraviMerge}. In this variant, we introduce a feedback force $\mathbf{F}_{\text{feedback}}^{(l)}$ pointing from the spacecraft toward the live normalized activation at layer $l-1$:
\vspace{-3pt}
\begin{equation}
    \mathbf{F}_{\text{feedback}}^{(l)} = \eta_{\text{feedback}} \left( \frac{\mathbf{h}^{(l-1)}}{\left\| \mathbf{h}^{(l-1)} \right\|_2} - \mathbf{h}_{\text{sc}}^{(l-1)} \right)
    \label{eq:feedback_force}
\end{equation}
\vspace{-3pt}
where $\eta_{\text{feedback}} \ge 0$ is a coupling coefficient. When $\eta_{\text{feedback}} > 0$, the net acceleration experiences a correction pull that couples the auxiliary spacecraft's trajectory directly to the actual propagating representations. Under Newton's second law, the total acceleration is updated as $\mathbf{a}^{(l)} = \sum_{k=1}^K \mathbf{F}_k^{(l)} + \mathbf{F}_{\text{feedback}}^{(l)}$, establishing a fully closed-loop dynamical system.

\textbf{Temporal State Carryover for Sequential Streams: } To support continuous, non-stationary task streams where queries arrive sequentially over time (resolving the statelessness critique), we can carry over the spacecraft's physical velocity vector $\mathbf{v}^{(L)}$ at the exit of one query to initialize the velocity vector of the next incoming query:
\vspace{-3pt}
\begin{equation}
    \mathbf{v}^{(3)}_{\text{next}} = \lambda_{\text{temporal}} \mathbf{v}^{(L)}_{\text{prev}}
    \label{eq:temporal_carryover}
\end{equation}
\vspace{-3pt}
where $\lambda_{\text{temporal}} \in [0, 1]$ is a temporal decay parameter. This introduces inter-query physical momentum, smoothing out transitions between different tasks in a continuous serving stream.

---

\vspace{-5pt}
\subsection{Control-Theoretic Foundation of Second-Order Representation Smoothing}
\label{sec:control_theoretic}
To address the concern regarding the heuristic nature of our physical analogy, we provide a control-theoretic justification for why second-order Newtonian dynamics are optimal for smoothing ensembling representation trajectories. Let activation-routing be modeled as a closed-loop dynamical system where ensembling weights update activations, altering the next layer's inputs, which feed back into the controller.

In control theory, first-order smoothers (EMA and ChemMerge) act as standard first-order low-pass filters with a transfer function $H_1(s) = 1 / (\tau s + 1)$. In a closed feedback loop, this first-order kinetics introduce a severe \textbf{phase lag} (delay) between the activation shift and the controller's response, causing the ensembling weights to overshoot the target centroid and trigger large-scale oscillations that destabilize the representation trajectory. This explains why EMA and ChemMerge suffer from massive accuracy penalties and high routing jitter.

In contrast, GraviMerge's second-order physical dynamics introduce virtual mass $m = 1.0$ and viscous drag $c = 1 - \gamma_{\text{drag}}$, forming a second-order spring-mass-damper system: $m \ddot{\mathbf{x}} + c \dot{\mathbf{x}} + k \mathbf{x} = \mathbf{F}$. This transfer function $H_2(s) = 1 / (m s^2 + c s + k)$ decays at $-40$ dB/decade at high frequencies (compared to $-20$ dB/decade for first-order filters), acting as a highly active low-pass filter that dampens high-frequency noise to near-zero ($0.00190$ MAD). Furthermore, unlike passive filters, GraviMerge is \textit{force-driven}. Our active Newtonian force $\mathbf{F}_k^{(l)}$ pulls the probe proactively toward the target centroid, while the velocity vector stores momentum, enabling the probe to converge smoothly without lag-induced task alignment delays. This control-theoretic proof demonstrates that modeling ensembling as a second-order physical system is a mathematically necessary framework to break the lag-accuracy barrier (see Section~\ref{sec:appendix_control_theoretic} in the Appendix for the full, detailed formulation).

---

\vspace{-5pt}
\subsection{Gravitational Influence Blending (GIB)}
Instead of using artificial heuristics to decay or threshold routing weights, the ensembling weight $\alpha_k^{(l)}$ for each specialist adapter is derived directly from the physical fields. Specifically, the blending coefficient $\alpha_k^{(l)}$ is proportional to the relative magnitude of the softened gravitational force exerted by star $k$ on the spacecraft probe:
\vspace{-3pt}
\begin{equation}
    \alpha_k^{(l)} = \frac{\left\| \mathbf{F}_k^{(l)} \right\|_2}{\sum_{j=1}^K \left\| \mathbf{F}_j^{(l)} \right\|_2} = \frac{M_k / (r_{k, l-1}^2 + \epsilon^2)}{\sum_{j=1}^K M_j / (r_{j, l-1}^2 + \epsilon^2)}
    \label{eq:gib}
\end{equation}
\vspace{-3pt}

By defining ensembling weights as continuous functions of the spacecraft's physical coordinates, we guarantee that $\alpha_k^{(l)}$ changes smoothly and differentiably from layer to layer. 

Crucially, the actual neural network activation vector $\mathbf{h}^{(l)}$ is updated using the standard ensembling formula (identical to the baselines), driven by these smooth physical weights $\alpha_k^{(l)}$:
\vspace{-3pt}
\begin{equation}
    \tilde{\mathbf{h}}^{(l)} = \mathbf{h}^{(l-1)} + \gamma_{\text{backbone}} \left( \sum_{k=1}^K \alpha_k^{(l)} \boldsymbol{\mu}_k^{(3)} - \mathbf{h}^{(l-1)} \right)
    \label{eq:activation_blend}
\end{equation}
\vspace{-3pt}
\begin{equation}
    \mathbf{h}^{(l)} = \frac{\tilde{\mathbf{h}}^{(l)}}{\left\| \tilde{\mathbf{h}}^{(l)} \right\|_2}
    \label{eq:activation_normalization}
\end{equation}
\vspace{-3pt}
where $\gamma_{\text{backbone}} \in [0, 1]$ is the backbone adapter update scale (calibrated to $0.3$). This formulation fully resolves the representational mismatch: the neural activations $\mathbf{h}^{(l)}$ propagate strictly through standard model ensembling operations, while the auxiliary spacecraft $\mathbf{h}_{\text{sc}}^{(l)}$ acts as an elegant stateful router on the side. This provides the best of both worlds: absolute mathematical coherence and perfect, jitter-free representational stability.
